% Transcription — fonds Alexandre Grothendieck, shelfmark 25, batch 7 (pages 121-140).
% Established from the facsimile https://grothendieck.umontpellier.fr/25.pdf, read on the
% private mirror archives/raw/25.pdf, per the skill transcribe-grothendieck.
% Pass: Opus 5.5 (claude-opus-5-5), 2026-10-03 — first pass, unchecked against the pages by a human.
%
% Limited pass for issue #46. Folder 25 (« Nouveau EGA V ») is covered by an edition made
% elsewhere (the EGA V prenotes), so only what that edition does not carry is transcribed:
% the typed letter to Dieudonné of pp. 134-135, dated in his hand 29.9.1965 and signed. The
% other pages of the batch (pp. 121-133 and 136-140: typescript of §§ IV 19-21 and the
% contents of « IV 21 », with his ink) are left out, which is why the page numbers jump.
%
% The letter is typed, with typed corrections (x-ed letters, a typed insertion between the
% lines) and his ink: the date at the head, the strike-through of the paragraph « Pour
% par.20 … » on p. 135, the signature. The archivists number the leaves 134-135; the typist's
% (or his) pencil numbers them 37-38, the letter having been filed among the notes.
\documentclass[11pt,a4paper]{article}
\input{../preamble/grothendieck}

\folder{25}
\batch{7}
\pages{121}{140}
\foldertitle{Nouveau EGA V : copies de tapuscrits annotés (s.d.).}
\dating{[à partir de 1967-1987]}
\watermark{Édition de démonstration}

\begin{document}

\page{134}
\note{lettre dactylographiée à Jean Dieudonné, seule partie du lot transcrite ici ; les pages 121 à 133 et 136 à 140 sont le tapuscrit des §§ IV 19 à 21, couvert par une autre édition}

\add{29.9.1965}

Cher Dieudonné,

Merci de ta lettre du 24 et pour la table des matières des par.
16 à 19. Je serais content de recevoir à l'occasion la table des
matières provisoire des par.20 et 21 ; d'accord pour les joindre au
fascicule 4 du Chap IV. Mais comment vaux-tu \struck{les} \add{mon ancien par.20} subdiviser, et
quels seront les titres des deux morceaux \uncertain{2}\note{sic, pour « ? »} Comme je commence à me
perdre dans le plan, et qu'il est \uncertain{paffois}\note{sic} commode de pouvoir référer
sans trop déconner à un nº de paragraphe, je te donne ici ce qui me
semble être le plan actuel, dis-moi si tu es d'accord :

20. ???

21. ???

22. Systêmes linéaires, compléments sur le groupe de Picard.

23. Grassmaniennes.

24. Formes lisses, singularités quadratiques ordinaires.

25. Sections hyperplanes et bordel.

26. Résultant et discriminant.

27. Extensions infinitésimales.

Le 25. risque d'ailleurs d'être fort long, et je te vois déjà
vouloir le subdiviser en deux ! Pourtant, $27=3^{3}$ est un bien joli
nombre !

Il n'est pas question que je publie l'ex-Appendice au par.18
sous mon nom ; ta rédaction n'a à peu près plus rien de commun avec
les vagues notes manuscriptes que je t'avais passées, si même je t'en
ai jamais passé, et ne me suis borné à te dire : il n'y a qu'à faire

\page{135}
pareil que pour les anneaux complets … Il serait d'autre part dommage
que ton travail de mise au point soit perdu pour les éventuels
utilisateurs (il finit toujours par s'en trouver ..). C'est pourquoi
je te demande de bien vouloir reconsidérer la question d'en faire un
« joint \uncertain{papar} »\note{sic, pour « paper »}.

\struck{Pour par.20, 10.9.1., il faut bien entendu utiliser le
fait que l'ensemble \struck{\ill{}} des points de $Z_{\lambda}$ en lesquels $F_{\lambda}$ \struck{\ill{}}
restreint à la fibre est de prof $\uncertain{\geq} n$ donné, est \emph{constructible}
(on a même dû prouver au par.12 qu'il est ouvert, avec les hypothèses
de platitude et de présentation finie qu'on a faites). Comme son image
inverse dans $Z$ est tout, c'est que c'est déjà tout un peu plus loin
que $\lambda$. C'est vraiment toujours le même argument qui revient !}
\note{paragraphe barré d'un trait oblique à l'encre}

Bien à toi

\add{A Grothendieck}
\note{signature autographe}

\end{document}
